\documentclass[11pt,a4paper]{article}

\usepackage[T1]{fontenc}
\usepackage[utf8]{inputenc}
\usepackage[french]{babel}
\usepackage{lmodern}
\usepackage{geometry}
\usepackage{graphicx}
\usepackage{hyperref}
\usepackage{microtype}
\usepackage{tabularx}

\geometry{margin=2.5cm}
\setlength{\emergencystretch}{2em}
\newcommand{\qbar}[1]{\rule{#1mm}{1.8ex}}
\newcolumntype{L}[1]{>{\raggedright\arraybackslash}p{#1}}
\newcolumntype{Y}{>{\raggedright\arraybackslash}X}
\hypersetup{
  colorlinks=true,
  linkcolor=black,
  citecolor=black,
  urlcolor=blue,
  pdftitle={État de l'art sur les LLMs et l'environnement},
  pdfauthor={Codex}
}

\title{État de l'art sur les grands modèles de langage et l'environnement}
\author{Document de synthèse académique}
\date{10 mars 2026}

\begin{document}

\maketitle

\begin{abstract}
Les grands modèles de langage (LLMs) sont devenus l'objet principal du débat sur l'impact environnemental de l'IA, mais la littérature reste difficile à comparer car elle mélange souvent des unités hétérogènes: entraînement, prompt, page générée, datacenter ou projection systémique. Cette version de l'état de l'art retient uniquement des données chiffrées rattachées à des sources primaires ou quasi primaires identifiables à la page, au tableau ou à la section. Elle montre trois faits robustes. Premièrement, l'entraînement d'un LLM frontier peut atteindre plusieurs milliers de tonnes de CO$_2$e, alors que certains cas open documentés restent dans des ordres de grandeur bien plus faibles. Deuxièmement, l'inférence unitaire peut être faible à l'échelle d'un prompt, mais devient stratégique à mesure que les usages se massivent. Troisièmement, l'évaluation environnementale ne peut plus se limiter au seul carbone de l'entraînement: eau, infrastructure, énergie de service et cycle de vie du matériel deviennent des dimensions centrales.
\end{abstract}

\textbf{Mots-clés:} LLMs, grands modèles de langage, environnement, CO$_2$, datacenters, Green AI.

\section{Introduction}
Le problème environnemental posé par les LLMs ne se réduit ni à un seul chiffre de CO$_2$ ni au seul événement d'entraînement. Il engage au moins quatre niveaux d'analyse: le coût de création des modèles, le coût unitaire d'inférence, l'infrastructure de datacenter nécessaire pour les servir, et enfin les usages concrets auxquels ces modèles sont affectés \cite{rillig2023,ren2024,larosa2025}. La littérature a longtemps été dominée par des estimations d'entraînement; depuis 2024-2025, elle s'élargit vers des mesures d'inférence, des analyses de cycle de vie et des travaux sur l'eau.

Cette section adopte une règle simple: ne retenir que les valeurs rattachées à une source identifiable dans le texte, un tableau ou une section nommée. Cela conduit à privilégier quelques points d'ancrage plutôt qu'une accumulation d'estimations hétérogènes. Le résultat est moins exhaustif en apparence, mais plus défendable scientifiquement.

\section{Infrastructure, électricité et intensification des usages}
Les LLMs doivent d'abord être replacés dans la dynamique générale des datacenters. L'IEA estime que les datacenters ont consommé environ 415 TWh d'électricité en 2024, soit environ 1,5\% de l'électricité mondiale, et projette environ 945 TWh en 2030 dans son scénario central \cite{iea2025}. Sur la même page, l'agence indique une croissance d'environ 15\% par an entre 2024 et 2030, précise que les serveurs accélérés tirés par l'IA représentent presque la moitié de l'augmentation nette de la demande électrique des datacenters et souligne que les Etats-Unis et la Chine concentreraient à eux deux près de 80\% de cette croissance \cite{iea2025}. Ces ordres de grandeur ne mesurent pas les LLMs seuls, mais ils fixent le cadre systémique de leur diffusion.

Le cas américain est encore plus documenté. Berkeley Lab rapporte que la consommation électrique des datacenters américains est passée de 58 TWh en 2014 à 176 TWh en 2023, avec une projection de 325 à 580 TWh en 2028; la part des datacenters dans la consommation nationale passerait alors d'environ 4,4\% en 2023 à 6,7\%--12\% en 2028 \cite{lbl2025}. L'article de synthèse du rapport souligne en outre que la demande électrique des datacenters a plus que doublé entre 2017 et 2023, largement du fait des serveurs IA \cite{lbl2025}.

EPRI ajoute un cadrage utile sur les usages. D'après le résumé de son rapport, les applications d'IA représenteraient déjà 10\%--20\% de l'électricité des datacenters, une requête ChatGPT serait de l'ordre de 2,9 Wh et les quatre principaux hyperscalers avaient plus que doublé leur consommation électrique entre 2017 et 2021 \cite{epri2024}. Ces deux derniers chiffres sont des estimations de rapport et non des mesures instrumentées comparables à celles de Google; ils sont donc utiles comme repère macro, pas comme base de comparaison physique fine.

\begin{table}[htbp]
\centering
\begin{tabularx}{\textwidth}{|L{3.4cm}|L{2.4cm}|Y|Y|}
\hline
Objet & Valeur & Référence & Localisation dans la source \\
\hline
Datacenters mondiaux & 415 TWh en 2024 & IEA (2025) \cite{iea2025} & Page web ``Energy demand from AI'', paragraphe ``The outlook for energy demand from data centres'' \\
\hline
Datacenters mondiaux & 945 TWh en 2030 & IEA (2025) \cite{iea2025} & Même paragraphe, scénario central \\
\hline
Croissance mondiale & $\sim$15\% par an, 2024--2030 & IEA (2025) \cite{iea2025} & Même paragraphe \\
\hline
Datacenters États-Unis & 58 TWh en 2014; 176 TWh en 2023; 325--580 TWh en 2028 & Berkeley Lab (2025) \cite{lbl2025} & Résumé du rapport, puce ``The total data center electricity usage climbed...'' \\
\hline
Part des datacenters dans l'électricité US & 4,4\% en 2023; 6,7\%--12\% en 2028 & Berkeley Lab (2025) \cite{lbl2025} & Résumé du rapport, puce précédente \\
\hline
Part de l'IA dans l'électricité des datacenters & 10\%--20\% & EPRI (2024) \cite{epri2024} & Résumé du rapport relayé par APPA, paragraphe central \\
\hline
Requête conversationnelle & $\sim$2,9 Wh par requête ChatGPT & EPRI (2024) \cite{epri2024} & Même paragraphe; chiffre de rapport cité comme estimation \\
\hline
\end{tabularx}
\caption{Repères macro sur l'électricité, les datacenters et les usages de l'IA.}
\label{tab:macro}
\end{table}

\section{Entraînement: du précédent historique aux LLMs frontier}
Le point de départ historique reste Strubell et al., qui montraient déjà avant la vague des LLMs génératifs que le coût climatique du NLP moderne pouvait devenir très élevé. Leur Table~1 rapporte 192 lb CO$_2$e pour un Transformer ``big'', 78\,468 lb CO$_2$e pour un pipeline NLP avec tuning et 626\,155 lb CO$_2$e pour une recherche d'architecture neuronale; leur Table~3 détaille ensuite les estimations de coût d'entraînement \cite{strubell2019}. Ce travail ne porte pas sur des LLMs frontier contemporains, mais il fixe la trajectoire historique du problème.

Le premier ancrage robuste dans l'ère des très grands modèles est BLOOM. Dans le résumé de leur article JMLR, Luccioni et al. estiment que l'entraînement final de BLOOM 176B a émis environ 24,7 tCO$_2$e en ne retenant que la consommation dynamique, et 50,5 tCO$_2$e en incluant un périmètre plus complet \cite{luccioni2023}. Leur Table~3 ventile ce total en 24,69 t pour la consommation dynamique, 14,6 t pour la consommation idle et 11,2 t pour les émissions incorporées, soit 50,5 tCO$_2$e au total \cite{luccioni2023}. Le point fort de ce papier est de montrer que le résultat varie fortement selon la frontière système retenue.

Le second ancrage est la documentation de Meta sur Llama 3.1. Dans la section ``Hardware and Software'' de la model card, Meta indique un total cumulé de 39,3 millions de GPU-hours sur H100-80GB pour la famille Llama 3.1, dont 30,84 millions pour le 405B; la même table attribue 8\,930 tCO$_2$e location-based au seul Llama~3.1~405B et 11\,390 tCO$_2$e à l'ensemble de la famille. La section ``Training Data'' ajoute que le pretraining a porté sur environ 15 trillions de tokens \cite{meta2024llama}. Méthodologiquement, cette source est très importante car elle documente à la fois l'échelle de calcul et les émissions, mais elle reste une model card plutôt qu'un article revu par les pairs.

Enfin, Morrison et al. déplacent l'analyse au-delà du run final. L'abstract de leur papier ICLR 2025 indique que la série de modèles étudiée, allant de 20 millions à 13 milliards de paramètres actifs et jusqu'à 5,6 trillions de tokens, a relâché 493 tonnes de carbone et consommé 2,769 millions de litres d'eau une fois comptés la fabrication du hardware, le développement et les runs finaux; les auteurs ajoutent que le développement représente à lui seul environ 50\% de l'impact de l'entraînement \cite{morrison2025}. Leur section de résultats rappelle aussi que la puissance observée fluctue entre environ 15\% et 85\% de la puissance maximale du matériel, ce qui complique les raisonnements fondés sur une puissance nominale constante \cite{morrison2025}.

\begin{table}[htbp]
\centering
\begin{tabularx}{\textwidth}{|L{3.3cm}|L{2.3cm}|L{2.6cm}|Y|}
\hline
Source & Métrique & Valeur & Localisation dans la source \\
\hline
Strubell et al. (2019) \cite{strubell2019} & CO$_2$e historique NLP & 192 lb; 78\,468 lb; 626\,155 lb & Table~1, p.~1 du PDF; Table~3, p.~4 du PDF \\
\hline
Luccioni et al. (2023) \cite{luccioni2023} & BLOOM 176B, entraînement final & 24,7 tCO$_2$e & Abstract, p.~1 du PDF \\
\hline
Luccioni et al. (2023) \cite{luccioni2023} & BLOOM 176B, périmètre étendu & 50,5 tCO$_2$e & Abstract, p.~1; Table~3, p.~7 du PDF \\
\hline
Meta Llama 3.1 (2024) \cite{meta2024llama} & Llama 3.1 405B, émissions location-based & 8\,930 tCO$_2$e & Section ``Hardware and Software'', tableau ``Training Location-Based Greenhouse Gas Emissions'' \\
\hline
Meta Llama 3.1 (2024) \cite{meta2024llama} & Llama 3.1 405B, calcul & 30,84M GPU-hours & Même tableau \\
\hline
Meta Llama 3.1 (2024) \cite{meta2024llama} & Corpus de pretraining & $\sim$15T tokens & Section ``Training Data'' \\
\hline
Morrison et al. (2025) \cite{morrison2025} & Série de modèles, carbone total & 493 tCO$_2$e & Abstract, p.~1 du PDF \\
\hline
Morrison et al. (2025) \cite{morrison2025} & Série de modèles, eau totale & 2,769 millions de litres & Abstract, p.~1; discussion quantitative p.~7 du PDF \\
\hline
Morrison et al. (2025) \cite{morrison2025} & Part du développement & $\sim$50\% de l'impact de l'entraînement & Abstract, p.~1 du PDF \\
\hline
\end{tabularx}
\caption{Valeurs d'entraînement retenues pour l'état de l'art.}
\label{tab:training}
\end{table}

Deux conclusions fortes se dégagent. D'une part, les ordres de grandeur vont de quelques dizaines de tonnes à plusieurs milliers de tonnes de CO$_2$e selon le modèle et le périmètre. D'autre part, les papiers les plus solides montrent que le ``vrai'' coût ne se limite ni à la seule énergie dynamique des accélérateurs ni au seul run final.

\section{Inférence: de l'estimation par requête à la mesure en production}
L'inférence est devenue la frontière centrale de la littérature récente. Le changement majeur depuis 2024 tient à l'apparition de mesures plus fines de l'énergie et de l'eau par prompt. Elsworth et al. proposent à ce titre la source la plus importante actuellement disponible pour un service en production. Leur abstract rapporte qu'un prompt texte médian dans Gemini Apps consomme 0,24 Wh, que les efforts d'efficacité logicielle et d'approvisionnement électrique bas carbone ont permis une réduction de 33x de l'énergie par prompt et de 44x de l'empreinte carbone en un an, et que le prompt médian consomme 0,26 mL d'eau \cite{elsworth2025}. Le Tableau~2 du papier donne ensuite 0,24 Wh/prompt, 0,03 gCO$_2$e/prompt et 0,26 mL/prompt pour l'approche de mesure ``comprehensive'' \cite{elsworth2025}. Le même passage ventile l'énergie: 0,14 Wh pour les accélérateurs actifs, 0,06 Wh pour CPU+DRAM, 0,02 Wh pour les machines idle provisionnées et 0,02 Wh pour l'overhead de datacenter \cite{elsworth2025}. Cette décomposition est importante parce qu'elle montre qu'une comparaison limitée au seul GPU sous-estime l'énergie réelle de service.

Ren et al. proposent un autre type d'unité: la génération d'une page de 500 mots. Dans la section ``Results'', ils estiment qu'un Llama-3-70B déployé dans un datacenter américain performant consomme 0,0195 kWh par page, émet 15 gCO$_2$ et consomme 0,14 litre d'eau une fois l'impact incorporé pris en compte \cite{ren2024}. Pour Gemma-2B-it, les mêmes auteurs obtiennent 0,00024 kWh, 0,18 gCO$_2$ et 0,0017 litre d'eau par page \cite{ren2024}. Ces résultats sont utiles pour discuter les unités de comparaison, mais ils reposent sur une analyse de cycle de vie comparative et non sur une mesure de produit en production comme chez Google.

L'estimation EPRI d'environ 2,9 Wh par requête ChatGPT reste pour sa part très citée \cite{epri2024}. Elle est utile comme ordre de grandeur de débat public, mais elle repose sur un rapport de cadrage plutôt que sur une publication détaillée de type benchmark ou mesure in situ. La tension apparente entre 2,9 Wh et 0,24 Wh ne doit donc pas être interprétée comme une contradiction simple: les frontières système, le produit considéré, la longueur de la requête et la méthode de mesure ne sont pas les mêmes.

\begin{table}[htbp]
\centering
\small
\begin{tabularx}{\textwidth}{|L{3.3cm}|L{2.3cm}|L{2.7cm}|Y|}
\hline
Source & Unité d'analyse & Valeurs retenues & Localisation dans la source \\
\hline
EPRI (2024) \cite{epri2024} & requête ChatGPT & $\sim$2,9 Wh/query & Résumé de rapport relayé par APPA, paragraphe central \\
\hline
Elsworth et al. (2025) \cite{elsworth2025} & prompt médian Gemini Apps & 0,24 Wh; 0,03 gCO$_2$e; 0,26 mL & Table~2, p.~7 du PDF \\
\hline
Elsworth et al. (2025) \cite{elsworth2025} & gains d'efficacité sur 12 mois & 33x moins d'énergie; 44x moins d'émissions & Abstract p.~1; Section~4.3, pp.~7--8 du PDF \\
\hline
Elsworth et al. (2025) \cite{elsworth2025} & décomposition énergétique du prompt & 0,14 Wh accélérateur; 0,06 Wh CPU+DRAM; 0,02 Wh idle; 0,02 Wh PUE & Section~4.2, p.~7 du PDF \\
\hline
Ren et al. (2024) \cite{ren2024} & page de 500 mots, Llama-3-70B & 0,0195 kWh; 15 gCO$_2$; 0,14 L & Section ``Results'', paragraphes sur Llama-3-70B \\
\hline
Ren et al. (2024) \cite{ren2024} & page de 500 mots, Gemma-2B-it & 0,00024 kWh; 0,18 gCO$_2$; 0,0017 L & Section ``Results'', paragraphes sur Gemma-2B-it \\
\hline
\end{tabularx}
\caption{Valeurs d'inférence retenues pour l'état de l'art.}
\label{tab:inference}
\end{table}

Le résultat clé est donc le suivant: l'inférence unitaire n'est pas nécessairement très élevée, mais elle est l'endroit où la massification des usages transforme un coût faible par interaction en enjeu d'infrastructure continu.

\section{Eau, cycle de vie et matériel}
L'eau est sortie du statut d'indicateur secondaire. L'article ``Making AI Less Thirsty'' montre que l'entraînement de GPT-3 dans des datacenters américains de Microsoft peut consommer au total 5,4 millions de litres d'eau, dont 700\,000 litres de consommation onsite scope~1; le même article ajoute qu'une bouteille de 500 mL correspond à environ 10 à 50 réponses de longueur moyenne selon le lieu et le moment de déploiement \cite{li2025water}. Ces chiffres sont méthodologiquement importants parce qu'ils distinguent eau directe et eau indirecte liée à l'électricité.

Morrison et al. confirment que l'eau doit être suivie sur l'ensemble du processus de création du modèle, pas seulement pendant le run final \cite{morrison2025}. Elsworth et al. montrent de leur côté qu'à l'échelle d'un prompt industriel très optimisé, l'eau peut tomber à 0,26 mL/prompt \cite{elsworth2025}. Pris ensemble, ces travaux indiquent que l'eau est encore plus dépendante que le carbone de la localisation, du mix électrique et de la technologie de refroidissement.

Sur la matérialité, la littérature spécifique aux LLMs reste plus lacunaire. Le résultat robuste n'est pas une comptabilité des terres rares par modèle, mais plutôt la montée des travaux de cycle de vie qui réintègrent fabrication du hardware, émissions incorporées et allocation du matériel au service rendu \cite{morrison2025}. C'est pourquoi il est prudent, à ce stade, de parler de carbone incorporé, de fabrication des accélérateurs et de renouvellement du hardware plutôt que de prétendre disposer d'une comptabilité minérale fine par LLM.

\begin{table}[htbp]
\centering
\begin{tabularx}{\textwidth}{|L{3.3cm}|L{2.3cm}|L{2.7cm}|Y|}
\hline
Source & Objet & Valeurs retenues & Localisation dans la source \\
\hline
Li et al. (2025) \cite{li2025water} & GPT-3, eau totale d'entraînement & 5,4 millions de litres, dont 700\,000 L onsite & Article CACM, paragraphe d'ouverture \\
\hline
Li et al. (2025) \cite{li2025water} & Chatbot génératif, eau d'usage & 500 mL pour 10--50 réponses moyennes & Même paragraphe d'ouverture \\
\hline
Morrison et al. (2025) \cite{morrison2025} & Série de modèles, eau totale & 2,769 millions de litres & Abstract, p.~1 du PDF \\
\hline
Elsworth et al. (2025) \cite{elsworth2025} & prompt médian Gemini Apps & 0,26 mL/prompt & Table~2, p.~7 du PDF \\
\hline
\end{tabularx}
\caption{Repères sur l'eau et le cycle de vie.}
\label{tab:water}
\end{table}

\section{Indicateurs cumulés annuels}
Les coûts unitaires par entraînement ou par inférence ne suffisent pas à apprécier l'importance systémique des LLMs. Pour une lecture macro, il faut disposer d'indicateurs cumulés annuels. La difficulté est que la littérature scientifique publie rarement une série mondiale annuelle ``LLMs only''. Les indicateurs cumulés les plus solides se situent donc à deux niveaux: d'une part, l'infrastructure générale des datacenters; d'autre part, les systèmes d'IA au sens large, qui incluent mais ne se limitent pas aux LLMs.

\subsection{Méthode de construction des courbes}
La logique retenue ici est volontairement simple et traçable.

\begin{itemize}
\item Lorsque la littérature publie directement une consommation annuelle en TWh, cette valeur est reprise telle quelle.
\item Lorsque la littérature publie une puissance moyenne ou un niveau de puissance pour une année donnée, la conversion en énergie annuelle est faite par la relation standard $TWh/an = GW \times 8{,}76$.
\item Les points observés et les points projetés sont distingués explicitement.
\item Aucune série annuelle n'est interpolée année par année quand la publication ne fournit que quelques points d'ancrage; on préfère montrer des ruptures de niveau plutôt qu'inventer une pseudo-précision.
\end{itemize}

Cette méthode permet de construire des courbes de synthèse tout en restant fidèle aux publications sources.

\subsection{Electricité annuelle des datacenters}
Le premier indicateur cumulé annuel disponible est la consommation électrique des datacenters. Le rapport de Shehabi et al. fournit pour les États-Unis trois points structurants: 58 TWh en 2014, 176 TWh en 2023, et 325--580 TWh en 2028 selon les scénarios \cite{shehabi2024report}. À l'échelle mondiale, l'IEA fournit 415 TWh en 2024 et 945 TWh en 2030 \cite{iea2025}. Ces valeurs ne sont pas propres aux LLMs, mais elles donnent le meilleur cadre annuel actuellement disponible pour situer leur diffusion.

\begin{table}[htbp]
\centering
\small
\begin{tabularx}{\textwidth}{|L{3.2cm}|L{2.2cm}|L{2.8cm}|Y|}
\hline
Périmètre & Année & Electricité annuelle & Source \\
\hline
Datacenters US & 2014 & 58 TWh & Shehabi et al. (2024) \cite{shehabi2024report} \\
\hline
Datacenters US & 2023 & 176 TWh & Shehabi et al. (2024) \cite{shehabi2024report} \\
\hline
Datacenters US & 2028 & 325--580 TWh & Shehabi et al. (2024), projection \cite{shehabi2024report} \\
\hline
Datacenters mondiaux & 2024 & 415 TWh & IEA (2025) \cite{iea2025} \\
\hline
Datacenters mondiaux & 2030 & 945 TWh & IEA (2025), projection \cite{iea2025} \\
\hline
\end{tabularx}
\caption{Points annuels cumulés utilisables pour une courbe de consommation électrique des datacenters.}
\label{tab:annual-datacenters}
\end{table}

\begin{figure}[htbp]
\centering
\begin{tabular}{l r l}
US 2014 & 58 & \qbar{10} \\
US 2023 & 176 & \qbar{30} \\
US 2028 bas & 325 & \qbar{55} \\
US 2028 haut & 580 & \qbar{98} \\
\end{tabular}
\caption{Repère visuel sur la hausse de la consommation électrique annuelle des datacenters aux États-Unis (TWh).}
\label{fig:annual-us-datacenters}
\end{figure}

\subsection{Puissance annuelle des systèmes d'IA et conversion en TWh}
Le second indicateur cumulé provient des travaux d'Alex de Vries-Gao, qui permettent d'approcher l'ampleur annuelle des systèmes d'IA. Dans Joule, l'auteur estime qu'en 2025 la puissance appelée par les systèmes d'IA sur l'ensemble de l'année pourrait se situer entre 5,3 et 9,4 GW, et que cette valeur pourrait atteindre 23 GW à la fin de 2025 \cite{devriesgao2025joule}. La version Patterns 2026 reprend également 9,4 GW à la fin de 2024 et jusqu'à 23 GW au cours de 2025 \cite{devriesgao2026patterns}.

En appliquant la conversion $GW \times 8{,}76$, on obtient les ordres de grandeur annuels suivants:

\begin{table}[htbp]
\centering
\small
\begin{tabularx}{\textwidth}{|L{4.2cm}|L{2.6cm}|L{2.8cm}|Y|}
\hline
Indicateur IA & Valeur de puissance & Equivalent annuel & Source \\
\hline
IA mondiale, moyenne basse 2025 & 5,3 GW & 46,4 TWh/an & de Vries-Gao (2025) \cite{devriesgao2025joule} \\
\hline
IA mondiale, moyenne haute 2025 & 9,4 GW & 82,3 TWh/an & de Vries-Gao (2025) \cite{devriesgao2025joule} \\
\hline
IA mondiale, niveau de fin 2025 annualisé & 23 GW & 201,5 TWh/an & de Vries-Gao (2025) \cite{devriesgao2025joule} \\
\hline
\end{tabularx}
\caption{Conversion en TWh annuels des puissances mondiales publiées pour les systèmes d'IA.}
\label{tab:annual-ai-power}
\end{table}

\begin{figure}[htbp]
\centering
\begin{tabular}{l r l}
IA 2025 basse & 46,4 & \qbar{8} \\
IA 2025 haute & 82,3 & \qbar{14} \\
IA fin 2025 annualisée & 201,5 & \qbar{34} \\
\end{tabular}
\caption{Ordres de grandeur annuels dérivés des puissances publiées pour les systèmes d'IA (TWh/an).}
\label{fig:annual-ai-power}
\end{figure}

\subsection{Portée et limites}
Ces courbes doivent être lues avec prudence. La courbe ``datacenters'' porte sur l'infrastructure totale, pas sur les seuls LLMs. La courbe ``IA mondiale'' issue de de Vries-Gao porte sur les systèmes d'IA au sens large, qui incluent les LLMs mais aussi d'autres charges. En revanche, ces deux familles d'indicateurs sont précisément celles qui permettent aujourd'hui de construire des synthèses annuelles défendables à partir de la littérature scientifique et technique indexée dans Google Scholar.

Le point important pour l'article est donc le suivant: il existe bien des indicateurs cumulés annuels permettant de montrer la dynamique d'augmentation des usages, de l'électricité et, indirectement, des pollutions. Mais au stade actuel de la littérature, ces courbes sont plus robustes au niveau `datacenters` et `AI systems` qu'au niveau ``LLMs mondiaux'' strictement défini. Cette distinction devra être explicitée dans le texte accompagnant les figures.

\section{Conclusion}
L'état actuel de la littérature permet de soutenir une position claire. L'entraînement des LLMs est maintenant suffisamment documenté pour établir des ordres de grandeur robustes, depuis BLOOM à quelques dizaines de tonnes de CO$_2$e jusqu'à Llama 3.1 405B à 8\,930 tCO$_2$e dans la documentation Meta \cite{luccioni2023,meta2024llama}. L'inférence est plus hétérogène, mais la mesure en production de Google montre qu'un prompt peut descendre à 0,24 Wh, 0,03 gCO$_2$e et 0,26 mL d'eau tout en restant un enjeu majeur à l'échelle des usages globaux \cite{elsworth2025}. Enfin, les travaux sur l'eau et le cycle de vie montrent que les évaluations mono-indicateur ne sont plus suffisantes \cite{li2025water,morrison2025}.

Pour une publication scientifique solide, la conséquence pratique est immédiate: toute valeur chiffrée doit être associée à une unité d'analyse, à une frontière système et à une localisation précise dans la source. Sans cette discipline, la revue de littérature redevient un assemblage de chiffres difficilement comparables. Avec elle, il devient possible de construire un état de l'art quantitatif qui ne soit pas du blabla, mais une base cumulative exploitable.

\bibliographystyle{plain}
\bibliography{references_ia_environnement}

\end{document}
